% source: https://github.com/pfradin/luadraw/discussions/137#discussioncomment-14892832 (pfradin, block 1)
\documentclass[12pt]{standalone}
\usepackage[svgnames]{xcolor}
\usepackage[3d]{luadraw}
\begin{document}
\begin{luadraw}{name=box23}
local x = 1.3
local win2d = {-5,10,-8,1.5}
local g = graph3d:new{window3d={0,7+2*x,0,10+2*x,-1,1}, margin={0,0,0,0},window={-5,5,-5,5},viewdir={30,70}, size={12,12}}
Hiddenlines = true; Hiddenlinestyle = "dashed" 

local draw_box = function(alpha, beta) -- angles of rotation
    local A, B, C, D, E, F, G = M(0,x,0), M(x,x,0), M(2*x,x,0), M(2*x,0,0), M(7+2*x,0,0), M(7+2*x,x,0), M(7+3*x,x,0)
    local N, M, L, K, J, I, H = table.unpack( sym3d({A,B,C,D,E,F,G}, {M(0,5+x,0),vecJ}) )
    A,B,M,N = table.unpack (rotate3d({A,B,M,N}, alpha, {C,vecJ}) ) --first we rotate A,B,M,N around (CL)
    A,N = table.unpack (rotate3d({A,N}, beta, {B,vecJ}) ) -- then we rotate A,N around (BM)
    D,E = table.unpack (rotate3d({D,E}, alpha, {C,-vecI}) ) -- and so on
    G,H = table.unpack (rotate3d({G,H}, alpha, {F,-vecJ}) )
    J,K = table.unpack (rotate3d({J,K}, alpha, {I,vecI}) )
    local facets = {{A,N,M,B},{B,M,L,C},{C,L,I,F},{D,C,F,E},{F,I,H,G},{L,K,J,I} }
    g:Dscene3d(
        g:addFacet(facets, {color="Orange",edge=true}))
end

-- first (top left)
g:Saveattr(); g:Viewport(-5,0,0,5); g:Coordsystem(table.unpack(win2d))
draw_box(0,0)
g:Restoreattr()

-- second (top right)
g:Saveattr(); g:Viewport(0,5,0,5); g:Coordsystem(table.unpack(win2d))
draw_box(20,45)
g:Restoreattr()

-- third (bottom left)
g:Saveattr(); g:Viewport(-5,0,-5,0); g:Coordsystem(table.unpack(win2d))
draw_box(45,90)
g:Restoreattr()

-- fourth (bottom right)
g:Saveattr(); g:Viewport(0,5,-5,0); g:Coordsystem(table.unpack(win2d))
draw_box(70,130)
g:Restoreattr()

g:Show() 
\end{luadraw}
\end{document}
